\subsection{Resultados de las mediciones}
\label{subsec:resultados-de-las-mediciones}

% ####################################################################################################################
% ####################################################################################################################

\subsubsection{Práctica 1 --- Ley de Ohm}

\todo{P1: tabla de pares $(I, V)$ con incertezas; gráfico $V$ vs.\ $I$ con barras de error y recta de ajuste; resultado final $R = m \pm \Delta m$ y ordenada $b \pm \Delta b$ (desarrollo en Apéndice~\ref{subsubsec:ajuste-cuadrados-minimos}).}

% ####################################################################################################################
% ####################################################################################################################

\subsubsection{Práctica 2 --- Mediciones directas e incertezas}

Las tres resistencias se midieron individualmente con el multímetro en modo $\Omega$, seleccionando en cada caso el rango más ajustado que contuviera el valor esperado.
Todas las mediciones son \textbf{directas}.
La incerteza se calculó según la Ecuación~\ref{eq:incerteza-multimetro} con los parámetros del rango utilizado; el desarrollo detallado se encuentra en el Apéndice~\ref{subsubsec:calculo-incertezas-resistencias}.

\begin{table}[H]
    \centering
    \caption{Especificaciones del multímetro UT58D (modo $\Omega$)}
    \label{tab:especificaciones-multimetro}
    \begin{tabular}{|c|c|c|c|}
        \hline
        \textbf{Rango} & \textbf{Resolución} & \textbf{Precisión} & \textbf{Corriente de prueba} \\ \hline
        200 $\Omega$   & 0{,}1 $\Omega$ & $\pm$0{,}7\% rdg $\pm$3 dgt & $<$ 0{,}7 mA \\ \hline
        2000 $\Omega$  & 1 $\Omega$     & $\pm$0{,}7\% rdg $\pm$1 dgt & $<$ 0{,}1 mA \\ \hline
        20 k$\Omega$   & 10 $\Omega$    & $\pm$0{,}7\% rdg $\pm$1 dgt & $<$ 30 $\mu$A \\ \hline
        200 k$\Omega$  & 100 $\Omega$   & $\pm$0{,}7\% rdg $\pm$1 dgt & $<$ 4 $\mu$A \\ \hline
        2000 k$\Omega$ & 1 k$\Omega$    & $\pm$1{,}0\% rdg $\pm$2 dgt & $<$ 0{,}4 $\mu$A \\ \hline
        20 M$\Omega$   & 10 k$\Omega$   & $\pm$2{,}0\% rdg $\pm$2 dgt & $<$ 40 nA \\ \hline
    \end{tabular}
\end{table}

\begin{table}[H]
    \centering
    \caption{Mediciones directas de resistencias (Práctica 2)}
    \label{tab:mediciones-resistencias}
    \begin{tabular}{|c|c|c|c|c|c|c|c|}
        \hline
        \textbf{R} & \textbf{Franjas} & \textbf{$R_{nom}$} & \textbf{Tolerancia} & \textbf{Rango} & \textbf{$R_{med}$} & \textbf{$\Delta R$} & \textbf{$R \pm \Delta R$} \\ \hline
        R1 & verde-azul-verde-dorado   & 5{,}6 M$\Omega$ & $\pm$5\% & 20 M$\Omega$   & 5{,}61 M$\Omega$ & 0{,}1 M$\Omega$ & $(5{,}6 \pm 0{,}1)$ M$\Omega$ \\ \hline
        R2 & verde-azul-amarillo-dorado & 560 k$\Omega$  & $\pm$5\% & 2000 k$\Omega$ & 562 k$\Omega$    & 8 k$\Omega$     & $(562 \pm 8)$ k$\Omega$ \\ \hline
        R3 & verde-azul-rojo-dorado    & 5{,}6 k$\Omega$ & $\pm$5\% & 20 k$\Omega$   & 5{,}50 k$\Omega$ & 0{,}05 k$\Omega$ & $(5{,}50 \pm 0{,}05)$ k$\Omega$ \\ \hline
    \end{tabular}
\end{table}

Los tres valores medidos se encuentran dentro del intervalo de tolerancia declarado ($\pm 5\%$):

\begin{itemize}
    \item R1: $[5{,}32;\, 5{,}88]$ M$\Omega$ $\ni$ 5{,}61 M$\Omega$ \checkmark
    \item R2: $[532;\, 588]$ k$\Omega$ $\ni$ 562 k$\Omega$ \checkmark
    \item R3: $[5{,}32;\, 5{,}88]$ k$\Omega$ $\ni$ 5{,}50 k$\Omega$ \checkmark
\end{itemize}

% ####################################################################################################################
% ####################################################################################################################

\subsubsection{Práctica 3 --- Influencia del aparato de medida}

\todo{P3: tabla de $V_{R_1}$, $V_{R_2}$, $V_{fuente}$ para cada par ($5{,}6\ \text{k}\Omega$; $560\ \text{k}\Omega$; $5{,}6\ \text{M}\Omega$); verificación de $\sum V = 0$; comparación ideal vs.\ real.}
